\documentclass[10pt,a4paper]{amsart}
\usepackage[margin=19mm]{geometry}
\usepackage{amsmath,amssymb,mathtools}
\usepackage[scr=rsfs,cal=euler]{mathalfa}
\usepackage{xfrac}
\usepackage[unicode,hidelinks,bookmarks=false]{hyperref}
\newcommand{\ie}{i.e.\ }
\newcommand{\cf}{cf.\ }
\newcommand{\resp}{respectively\ }
\newcommand{\Subsection}{\subsection}
\providecommand{\nobreakdash}{\nobreak\hbox{-}}
\setlength{\emergencystretch}{2em}
\multlinegap0pt
\usepackage{bm}
\usepackage{bbm}
\usepackage{dsfont}
\usepackage{xspace}
\usepackage{cancel}
\usepackage{mathtools}
\usepackage{amsthm}
\makeatletter
\@ifundefined{theorem}{\newtheorem{theorem}{Theorem}}{}
\@ifundefined{lemma}{\newtheorem{lemma}[theorem]{Lemma}}{}
\makeatother
\newcommand{\RB}{\mathbb{R}}
\title[Convergence Rates for Gradient Descent on the Edge of Stabil]{Convergence Rates for Gradient Descent on the Edge of Stability in Overparametrised Least Squares}
\author{Lachlan Ewen MacDonald, Hancheng Min, Leandro Palma, Salma Tarmoun, Ziqing Xu, Ren\'e Vidal}
\date{Source: arXiv:2510.17506v1 (CC BY 4.0)}
\begin{document}
\maketitle

\noindent\emph{Source and licence.} Convergence Rates for Gradient Descent on the Edge of Stability in Overparametrised Least Squares. Version arXiv:2510.17506v1 (CC BY 4.0), distributed under the Creative Commons Attribution 4.0 International licence, as recorded in the arXiv licence metadata for this exact version and shown on the abstract page https://arxiv.org/abs/2510.17506v1. Full rights notice: licenses/arxiv-2510.17506-CC-BY-4.0.txt.\par\smallskip

\noindent\textit{Editorial scope.} Counted unit: the Lemma whose source label is \texttt{lem:technical1} in the article (source0.tex lines 1814--1820, source label \texttt{lem:technical1}), delivered together with the article's own complete original proof of it (source0.tex lines 1822--1896, one \texttt{proof} environment of 75 lines). No mathematics is written, repaired or re-derived: statement and proof are the article's own text line for line. Required context retained uncounted: lem:invariant (lines 1792--1794). Every definition, display and label of those lines is retained unchanged. Omitted from this package and not redistributed: 1--1791, 1795--1813, 1821--1821, 1897--2252. Nothing inside a retained excerpt is omitted or altered: every retained body is delivered exactly as the article's lines are written, so no omission receipt of any kind accompanies this package. Citations: the retained excerpts carry no citation command at all, so no bibliography entry of the article is transcribed and no reference is redistributed. Reference adaptation: none was needed and none was made; every reference inside the delivered unit resolves inside it, so no unresolved reference or citation is produced. Printed slips kept literal: no printed word, symbol, hypothesis or number of the counted statement or of its proof was altered. Layout and class: the article's own class is \texttt{article} with packages including neurips\_2025, inputenc, fontenc, hyperref, url, booktabs, amsfonts, nicefrac, microtype, xcolor, graphicx, wrapfig, tcolorbox, amsmath; the wrapper is the lane's pinned \texttt{amsart} editorial wrapper at 10pt on A4, which restates the theorem-like environments the retained text needs and adds the article's own macro definitions that the retained text needs (\texttt{\textbackslash RB}), with extra wrapper packages (bm, bbm, dsfont, xspace, cancel, mathtools). The delivered numbering is the unit's own. Assets: no retained span contains a figure environment, a graphics-inclusion command, a PDF-page inclusion, a table or a TikZ picture; no image file is copied into this package, the package declares no asset dependency and no image file is redistributed with it. Licence: version arXiv:2510.17506v1 of this article is distributed under the Creative Commons Attribution 4.0 International licence, as recorded in the arXiv licence metadata for this exact version and shown on the abstract page https://arxiv.org/abs/2510.17506v1. A full rights notice is delivered at licenses/arxiv-2510.17506-CC-BY-4.0.txt. Characters: every retained excerpt is pure ASCII, so the delivered engine, XeLaTeX with its default Latin Modern text font, reports no missing character.
\begin{lemma}\label{lem:invariant}
    There exists $r>0$ and a function $u:\{x>0:|x^2-r|<r\}\rightarrow\RB$ whose graph $\Gamma_u:=\{(x,u(x)):x\in\text{dom}(u)\}$ is invariant under $T$, and for which $u(x) = \sqrt{\frac{b}{c-a}}x+O(x^2)$.
\end{lemma}

\begin{lemma}\label{lem:technical1}
    Let $u(x) = \sqrt{\frac{b}{c-a}}x+O(x^2)$ be the function from Lemma \ref{lem:invariant}, and define $\Delta(x,y):=y-u(x)$. Then $\widetilde{\Delta}(x,y):=\Delta(x,y)/x$ satisfies
    \begin{align}
        |\widetilde{\Delta}\circ T(x,y)|\leq |\widetilde{\Delta}(x,y)|\bigg(1-\frac{c-a}{2}y^2\bigg)
    \end{align}
    for all sufficiently small $x>0$, $y\geq 0$.
\end{lemma}

\begin{proof}
    We estimate $\frac{\widetilde{\Delta}\circ T}{\widetilde{\Delta}}$. Define
    \begin{align}
        D(x,y):=\frac{T_x(x,y)}{x},\qquad N(x,y):=\frac{\Delta\circ T(x,y)}{\Delta(x,y)},\qquad (x,y)\in\text{dom}(u)\times\RB_{\geq 0}.
    \end{align}
    Observe that both $D$ and $N$ are analytic on the domain of $u$. Observe moreover that
    \begin{align}
        \frac{\widetilde{\Delta}\circ T(x,y)}{\widetilde{\Delta}(x,y)} = \frac{\Delta\circ T(x,y)}{T_x(x,y)}\cdot \frac{x}{\Delta(x,y)} = \frac{N(x,y)}{D(x,y)} = 1-\frac{D(x,y)-N(x,y)}{D(x,y)},
    \end{align}
    so that estimation of $\frac{\widetilde{\Delta}\circ T}{\widetilde{\Delta}}$ is reduced to estimating $\frac{D-N}{D}$. We first estimate $D-N$. For this, observe that since $T(x,0) = (x,0)$, one has $D(x,0) - N(x,0) = 0$ for all $x\in\text{dom}(u)$. Thus, by analyticity of $D$ and $N$, there is a unique analytic function $H$ on $\text{dom}(u)\times\RB$ such that 
    \begin{align}\label{eq:Hdef}
        D(x,y) - N(x,y)= yH(x,y).
    \end{align}
    We compute $H(x,y)$ to second order. For notational convenience, set
    \begin{align}
        \kappa:=\sqrt{\frac{b}{c-a}}
    \end{align}
    so that $u(x) = \kappa x+O(x^2)$. Differentiate \eqref{eq:Hdef} with respect to $y$ and evaluate at $y=0$ to obtain
    \begin{align}
        H(x,0) = \partial_yD(x,0)-\partial_yN(x,0),
    \end{align}
    Observe that
    \begin{align}
        \partial_yD(x,0) = \partial_y(1-ay^2 + A(x,y)y^2)|_{y=0} = 0
    \end{align}
    while, from $N = (\Delta\circ T)/\Delta$, one has
    \begin{align}
        \partial_yN(x,0) &= \frac{\Delta\,\partial_y(\Delta\circ T)-(\Delta\circ T)\,\partial_y\Delta}{\Delta^2}(x,0).\label{eq:firstpartialN}
    \end{align}
    From $\Delta(x,y) = y-u(x)$, is easily seen that
    \begin{align}
        \Delta(x,0) = -u(x),\qquad \partial_y\Delta(x,0) = 1,\qquad \Delta\circ T(x,0) = -u(T_x(x,0)) = -u(x)
    \end{align}
    while
    \begin{align}
        \partial_y(\Delta\circ T)(x,0) = \partial_y T_y(x,0)-u'(x)\partial_yT_x(x,0) = 1+bx^2 + O(x^3)
    \end{align}
    so that \eqref{eq:firstpartialN} yields
    \begin{align}
        H(x,0) = -\partial_yN(x,0) = \frac{bx^2 + O(x^3)}{u(x)} = (c-a)\kappa x + O(x^2).\label{eq:Hfirst}
    \end{align}
    To compute the second order component of $H$, differentiate \eqref{eq:Hdef} twice with respect to $y$ and evaluate at $y=0$ to obtain
    \begin{align}
        \partial_yH(x,0) = \frac{1}{2}\big(\partial^2_yD(x,0)-\partial_y^2N(x,0)\big).\label{eq:partialH}
    \end{align}
    One sees that
    \begin{align}
        \partial^2_yD(x,0) = \partial^2_y(1-ay^2 + A(x,y)y^2)|_{y=0} = -2a+O(x),
    \end{align}
    while
    \begin{align}
        \partial^2_yN(x,0) = \bigg(\frac{\partial^2_y(\Delta\circ T)}{\Delta} - \frac{(\Delta\circ T)\,\partial^2_y\Delta}{\Delta^2} - 2\frac{\partial_y(\Delta\circ T)\,\partial_y\Delta}{\Delta^2} + 2\frac{(\Delta\circ T)\,(\partial_y\Delta)^2}{\Delta^3}\bigg)(x,0).
    \end{align}
    Since
    \begin{align}
        \partial^2_y(\Delta\circ T)(x,0) = \partial^2_yT_y(x,0) - u''(x)(\partial_yT_x(x,0))^2 - u'(x)\partial^2_yT_x(x,0) = 2a\kappa x+O(x^2)
    \end{align}
    one obtains
    \begin{align}
        \partial^2_yN(x,0) &= -\frac{2a\kappa x+O(x^2)}{\kappa x+O(x^2)}-2\frac{1+bx^2+O(x^3)}{\kappa^2x^2+O(x^3)} + 2\frac{1}{\kappa^2x^2 + O(x^3)}\\&=-2a-2b/\kappa^2 + O(x) = -2c+O(x).
    \end{align}
    Thus
    \begin{align}
        \partial_yH(x,0) = \frac{1}{2}(-2a+2c+O(x)) = (c-a)+O(x).\label{eq:Hsecond}
    \end{align}
    We deduce from \eqref{eq:Hfirst}, \eqref{eq:Hsecond} and \eqref{eq:Hdef} that
    \begin{align}
        D(x,y)-N(x,y) = (c-a)(y^2+\kappa yx) + O(yx^2,y^2x,y^3).
    \end{align}
    It follows that for all sufficiently small $x>0, y\geq0$, one has $D(x,y)-N(x,y)\geq \frac{(c-a)}{2}y^2$ and $1\geq D(x,y)>0$ so that
    \begin{align}
        \frac{\widetilde{\Delta}\circ T(x,y)}{\widetilde{\Delta}(x,y)} = 1-\frac{D(x,y)-N(x,y)}{D(x,y)}\leq 1-\frac{c-a}{2}y^2
    \end{align}
    as claimed.
\end{proof}

\end{document}
